% ╔═══════════════════════════════════════╗
% ║   CHAPITRE 1 : LOGIQUE               ║
% ╚═══════════════════════════════════════╝
\chapter{Logique et Raisonnements}

\section{Pourquoi la logique est fondamentale}

Les mathématiques reposent sur un langage précis là où la langue naturelle est ambiguë. Le mot <<~ou~>> en français peut être exclusif (<<~fromage ou dessert~>>) ou inclusif (<<~les as ou les c\oe urs~>>). La logique formelle lève cette ambiguïté.

\begin{intuitionbox}[Lien avec la cybersécurité]
En cybersécurité, la logique formelle est omniprésente : les règles de pare-feu, les politiques de contrôle d'accès, et même les requêtes SQL reposent sur des opérations booléennes. Comprendre \texttt{AND}, \texttt{OR}, \texttt{NOT} en logique, c'est comprendre les fondements de la sécurité informatique.
\end{intuitionbox}

\section{Assertions et connecteurs logiques}

\begin{defbox}[--- Assertion]
Une \textbf{assertion} (ou proposition logique) est un énoncé qui est soit \textbf{vrai}, soit \textbf{faux}, mais jamais les deux simultanément.
\end{defbox}

\subsection{Les connecteurs fondamentaux}

À partir de deux assertions $P$ et $Q$, on construit de nouvelles assertions :

\begin{defbox}[--- Connecteurs logiques]
\begin{itemize}[leftmargin=1.5cm]
  \item \textbf{Conjonction} : <<~$P$ et $Q$~>> est vraie si et seulement si $P$ est vraie \emph{et} $Q$ est vraie.
  \item \textbf{Disjonction} : <<~$P$ ou $Q$~>> est vraie si \emph{au moins} l'une des deux est vraie.
  \item \textbf{Négation} : <<~non $P$~>> est vraie si $P$ est fausse, et réciproquement.
\end{itemize}
\end{defbox}

\subsection{Tables de vérité}

\begin{center}
\begin{tikzpicture}[
  every node/.style={font=\small\sffamily},
  table/.style={
    matrix of nodes,
    row sep=-\pgflinewidth,
    column sep=-\pgflinewidth,
    nodes={draw, minimum width=1.2cm, minimum height=0.7cm, anchor=center},
    row 1/.style={nodes={fill=deepblue, text=white, font=\small\sffamily\bfseries}},
    column 1/.style={nodes={fill=lightblue}},
  }
]

% ET
\matrix[table] (et) at (0,0) {
  $P \wedge Q$ & V & F \\
  V & V & F \\
  F & F & F \\
};
\node[above=0.3cm, font=\sffamily\bfseries\color{deepblue}] at (et.north) {ET ($\wedge$)};

% OU
\matrix[table] (ou) at (5,0) {
  $P \vee Q$ & V & F \\
  V & V & V \\
  F & V & F \\
};
\node[above=0.3cm, font=\sffamily\bfseries\color{deepblue}] at (ou.north) {OU ($\vee$)};

% NON
\matrix[table] (non) at (10,0) {
  $P$ & non $P$ \\
  V & F \\
  F & V \\
};
\node[above=0.3cm, font=\sffamily\bfseries\color{deepblue}] at (non.north) {NON ($\neg$)};

\end{tikzpicture}
\end{center}

\subsection{L'implication}

\begin{defbox}[--- Implication]
L'implication $P \Rightarrow Q$ est définie comme l'assertion <<~$(\text{non } P) \text{ ou } Q$~>>. Autrement dit :

\medskip
$P \Rightarrow Q$ est \textbf{fausse} uniquement lorsque $P$ est vraie et $Q$ est fausse.

\medskip
En particulier, si $P$ est fausse, alors $P \Rightarrow Q$ est \textbf{toujours vraie}, quel que soit $Q$.
\end{defbox}

\begin{intuitionbox}[Comprendre le <<~faux implique tout~>>]
Ce point est contre-intuitif. L'astuce : l'implication $P \Rightarrow Q$ est une \emph{promesse}. Si je dis <<~s'il pleut, je prends un parapluie~>>, la promesse n'est brisée que dans un cas : il pleut et je n'ai pas de parapluie. S'il ne pleut pas, je n'ai rien promis, donc la promesse est respectée dans tous les cas. C'est exactement la ligne <<~F $\Rightarrow$ \dots = V~>> de la table de vérité.
\end{intuitionbox}

\subsection{L'équivalence}

\begin{defbox}[--- Équivalence]
L'équivalence $P \Leftrightarrow Q$ est l'assertion <<~$(P \Rightarrow Q)$ et $(Q \Rightarrow P)$~>>.

Elle est vraie lorsque $P$ et $Q$ ont la même valeur de vérité (toutes deux vraies ou toutes deux fausses).
\end{defbox}

\subsection{Lois fondamentales}

\begin{propbox}[--- Lois de De Morgan et distributivité]
\begin{enumerate}[leftmargin=1.5cm]
  \item $\text{non}(P \text{ et } Q) \;\Leftrightarrow\; (\text{non } P) \text{ ou } (\text{non } Q)$
  \item $\text{non}(P \text{ ou } Q) \;\Leftrightarrow\; (\text{non } P) \text{ et } (\text{non } Q)$
  \item $P \text{ et } (Q \text{ ou } R) \;\Leftrightarrow\; (P \text{ et } Q) \text{ ou } (P \text{ et } R)$
  \item $(P \Rightarrow Q) \;\Leftrightarrow\; (\text{non } Q \Rightarrow \text{non } P)$ \quad \textit{(contraposée)}
\end{enumerate}
\end{propbox}

\begin{proofbox}
Pour la loi de De Morgan (1), on vérifie par table de vérité que les deux colonnes sont identiques pour toutes les combinaisons de valeurs de $P$ et $Q$ :

\smallskip
\begin{center}
\begin{tabular}{cc|c|c}
$P$ & $Q$ & $\text{non}(P \wedge Q)$ & $(\neg P) \vee (\neg Q)$ \\
\hline
V & V & F & F \\
V & F & V & V \\
F & V & V & V \\
F & F & V & V \\
\end{tabular}
\end{center}
\smallskip

Les deux colonnes sont identiques, donc les assertions sont logiquement équivalentes. \hfill$\square$
\end{proofbox}

\begin{intuitionbox}[Pattern structurel des lois de De Morgan]
La règle génératrice est simple : la négation <<~traverse~>> un connecteur en le transformant en son dual :
\begin{itemize}[leftmargin=1cm]
  \item $\text{non}(\text{et}) \longrightarrow \text{ou}$
  \item $\text{non}(\text{ou}) \longrightarrow \text{et}$
\end{itemize}
C'est la même structure que $\overline{A \cap B} = \overline{A} \cup \overline{B}$ en théorie des ensembles.
\end{intuitionbox}


\section{Quantificateurs}

\begin{defbox}[--- Quantificateurs]
\begin{itemize}[leftmargin=1.5cm]
  \item \textbf{Universel} : $\forall x \in E,\; P(x)$ est vraie si $P(x)$ est vraie pour \emph{tout} élément $x$ de $E$.
  \item \textbf{Existentiel} : $\exists x \in E,\; P(x)$ est vraie s'il existe \emph{au moins un} $x$ dans $E$ tel que $P(x)$ soit vraie.
\end{itemize}
\end{defbox}

\subsection{Négation des quantificateurs}

\begin{propbox}[--- Règles de négation]
\[
  \text{non}\bigl(\forall x \in E,\; P(x)\bigr) \;\Leftrightarrow\; \exists x \in E,\; \text{non } P(x)
\]
\[
  \text{non}\bigl(\exists x \in E,\; P(x)\bigr) \;\Leftrightarrow\; \forall x \in E,\; \text{non } P(x)
\]
\end{propbox}

\begin{intuitionbox}[Algorithme de négation]
Pour nier une phrase logique, on applique mécaniquement :
\begin{enumerate}[leftmargin=1cm]
  \item $\forall \longleftrightarrow \exists$ (on échange les quantificateurs)
  \item On nie la proposition finale
  \item Les inégalités strictes deviennent larges et vice versa
\end{enumerate}
\textbf{Exemple :} La négation de $\forall \varepsilon > 0,\; \exists N \in \N,\; \forall n \geqslant N,\; |u_n - \ell| \leqslant \varepsilon$ est :
\[
  \exists \varepsilon > 0,\; \forall N \in \N,\; \exists n \geqslant N,\; |u_n - \ell| > \varepsilon
\]
C'est exactement <<~la suite ne converge pas vers $\ell$~>>.
\end{intuitionbox}

\subsection{L'ordre des quantificateurs compte}

\begin{exbox}[Permutation interdite]
Les deux assertions suivantes sont \textbf{différentes} :
\begin{align*}
  &\forall x \in \R,\; \exists y \in \R,\; x + y > 0 \quad \text{(\textbf{vraie} : prendre $y = |x| + 1$)} \\
  &\exists y \in \R,\; \forall x \in \R,\; x + y > 0 \quad \text{(\textbf{fausse} : aucun $y$ fixé ne convient pour tout $x$)}
\end{align*}
Dans la première, $y$ \emph{dépend} de $x$. Dans la seconde, $y$ doit être \emph{universel}.
\end{exbox}


\section{Méthodes de raisonnement}

\subsection{Raisonnement direct}

Pour montrer $P \Rightarrow Q$ : on suppose $P$ vraie et on montre $Q$.

\begin{exbox}
\textbf{Montrer que si $a, b \in \Q$ alors $a + b \in \Q$.}

On écrit $a = \frac{p}{q}$ et $b = \frac{p'}{q'}$ avec $p, p' \in \Z$ et $q, q' \in \N^*$. Alors :
\[
  a + b = \frac{pq' + qp'}{qq'} \in \Q
\]
car $pq' + qp' \in \Z$ et $qq' \in \N^*$. \hfill$\square$
\end{exbox}

\subsection{Raisonnement par contraposée}

On exploite l'équivalence : $(P \Rightarrow Q) \;\Leftrightarrow\; (\text{non } Q \Rightarrow \text{non } P)$.

\begin{exbox}
\textbf{Montrer que si $n^2$ est pair alors $n$ est pair.}

\textit{Contraposée :} on montre que si $n$ est impair alors $n^2$ est impair. Si $n = 2k+1$, alors $n^2 = 4k^2 + 4k + 1 = 2(2k^2 + 2k) + 1$, qui est bien impair. \hfill$\square$
\end{exbox}

\subsection{Raisonnement par l'absurde}

On suppose $P$ vraie et $Q$ fausse simultanément, puis on dérive une contradiction.

\begin{exbox}
\textbf{Montrer :} si $a, b \geqslant 0$ et $\frac{a}{1+b} = \frac{b}{1+a}$, alors $a = b$.

Supposons $a \neq b$. De $a(1+a) = b(1+b)$, on tire $a^2 - b^2 = b - a$, soit $(a-b)(a+b) = -(a-b)$. Comme $a \neq b$, on simplifie : $a + b = -1$. Contradiction car $a, b \geqslant 0$. \hfill$\square$
\end{exbox}

\subsection{Récurrence}

\begin{methodbox}[Raisonnement par récurrence]
Pour montrer $P(n)$ pour tout $n \geqslant n_0$ :
\begin{enumerate}[leftmargin=1.5cm]
  \item \textbf{Initialisation :} vérifier $P(n_0)$.
  \item \textbf{Hérédité :} fixer $n \geqslant n_0$, supposer $P(n)$ vraie (\textit{hypothèse de récurrence}), montrer $P(n+1)$.
  \item \textbf{Conclusion :} par le principe de récurrence, $P(n)$ est vraie pour tout $n \geqslant n_0$.
\end{enumerate}
\end{methodbox}

\begin{exbox}
\textbf{Montrer que $2^n \geqslant n$ pour tout $n \in \N$.}

\textit{Initialisation :} $2^0 = 1 \geqslant 0$. \checkmark

\textit{Hérédité :} supposons $2^n \geqslant n$. Alors $2^{n+1} = 2 \cdot 2^n \geqslant 2n \geqslant n + 1$ (car $n \geqslant 1$).

Pour $n = 0$ : $2^1 = 2 \geqslant 1$. \checkmark \hfill$\square$
\end{exbox}


\section{Synthèse personnelle --- Logique}

\begin{intuitionbox}[Connexions identifiées]
\textbf{1. La logique comme fondation unificatrice.} Tout le reste des mathématiques (suites, complexes, analyse, algèbre) s'appuie sur ce langage. La définition $\varepsilon$-$N$ de la limite d'une suite, la caractérisation d'un point fixe --- tout cela n'est que de la logique appliquée.

\textbf{2. Lien avec la programmation.} Les lois de De Morgan apparaissent directement dans l'optimisation de conditions booléennes en code. Le raisonnement par récurrence est l'analogue mathématique de la preuve de correction d'un algorithme récursif.

\textbf{3. Pattern de la contraposée.} Quand une implication directe semble difficile (car la conclusion donne peu de prise), essayer la contraposée. C'est un réflexe que j'ai intégré : <<~si la conclusion est plus simple à nier que l'hypothèse n'est à exploiter, prendre la contraposée~>>.
\end{intuitionbox}


